% ===== Document Preamble =====
\documentclass[11pt]{article}
\usepackage[margin=0.5in]{geometry}
\usepackage{amsmath,amssymb}
\usepackage{xcolor}
\usepackage{enumitem}
\usepackage{needspace}
\usepackage{comment}

% ===== Colors =====
\definecolor{AnsRed}{RGB}{170,0,0}
\definecolor{WrongBlue}{RGB}{0,0,160}
\definecolor{SkipGray}{RGB}{120,120,120}

% ===== Choice list styling (tight, consistent) =====
\setlist[enumerate]{%
  leftmargin=1.5em,
  itemsep=0.25em,
  topsep=0.25em,
  label=\textbf{\Alph*.}
}

% ===== Style Wrappers =====
\newcommand{\QuestionStyle}[1]{{\large\bfseries\boldmath #1}}
\newcommand{\AnswerStyle}[1]{{\large\bfseries\boldmath #1}}
\newcommand{\ReasonStyle}[1]{{\normalsize #1}}

% ===== Show / Hide Reasoning Toggle =====
\newif\ifshowreason
\showreasontrue      % Instructor version
% \showreasonfalse   % Student version

% ===== Choice Item Macros (Answer + Reasoning) =====
% Usage:
%   \Correct{Answer}{Reasoning}
%   \Wrong{Answer}{Reasoning}

\newcommand{\Correct}[2]{%
  \item \AnswerStyle{\textcolor{AnsRed}{#1}}%
  \ifshowreason\\[-0.2em]
    \ReasonStyle{\textcolor{AnsRed}{#2}}%
  \fi
}

\newcommand{\Wrong}[2]{%
  \item \AnswerStyle{\textcolor{WrongBlue}{#1}}%
  \ifshowreason\\[-0.2em]
    \ReasonStyle{\textcolor{WrongBlue}{#2}}%
  \fi
}

% ===== Multiple-Choice Question Block Macro =====
% Usage: \QBlock{Q1}{stem...}{<items>}{answer key}

\newcommand{\QBlock}[4]{%
  \Needspace{8\baselineskip}%
  \noindent
  \begin{minipage}{\linewidth}
    \QuestionStyle{#1.~#2}

    \vspace{0.25em}
    \begin{enumerate}
      #3
    \end{enumerate}

    \vspace{0.1em}
    \noindent\textcolor{AnsRed}{\textit{\textbf{Answer:}~#4}}

    \vspace{0.25em}
    \hrule height 0.4pt
  \end{minipage}
  \vspace{0.75em}
}

% ===== Open-Ended Question Block Macro =====
% Usage: \OEBlock{Q1}{stem...}{answer text}

\newcommand{\OEBlock}[3]{%
  \Needspace{5\baselineskip}%
  \noindent
  \begin{minipage}{\linewidth}
    \QuestionStyle{#1.~#2}

    \vspace{0.25em}
    \noindent\textcolor{AnsRed}{\textit{\textbf{Answer:}~#3}}

    \vspace{0.25em}
    \hrule height 0.4pt
  \end{minipage}
  \vspace{0.75em}
}

% ===== Multi-Part Question Block Macro =====
% Usage: \MPBlock{Q1}{stem...}{<parts>}

\newcommand{\MPBlock}[3]{%
  \Needspace{8\baselineskip}%
  \noindent
  \begin{minipage}{\linewidth}
    \QuestionStyle{#1.~#2}

    \vspace{0.25em}
    #3

    \vspace{0.25em}
    \hrule height 0.4pt
  \end{minipage}
  \vspace{0.75em}
}

% ===== MC Sub-Part Macro =====
% Usage: \MCPart{a}{stem...}{<items>}{answer key}

\newcommand{\MCPart}[4]{%
  \par\vspace{0.35em}
  \noindent\textbf{(#1)}~#2
  \begin{enumerate}
    #3
  \end{enumerate}
  \vspace{0.1em}
  \noindent\textit{\textbf{Answer:}~#4}
  \par
}

% ===== Open-Ended Sub-Part Macro =====
% Usage: \OEPart{a}{stem...}{answer text}

\newcommand{\OEPart}[3]{%
  \par\vspace{0.35em}
  \noindent\textbf{(#1)}~#2
  \par\vspace{0.15em}
  \noindent\textit{\textbf{Answer:}~#3}
  \par
}

% ===== Skipped Question Note Macro =====
% Usage: \SkipBlock{Q1}{reason}

\newcommand{\SkipBlock}[2]{%
  \noindent\textcolor{SkipGray}{\textit{[#1 omitted: #2]}}%
  \vspace{0.5em}\par
}

% ===== Document Start =====

\begin{document}

\begin{center}
{\huge \textbf{Physics Questions}}\\[0.25em]
{\Large Kinematics and Vectors}\\[0.25em]
{\large Questions generated below. Multiple-choice answers in RED; incorrect answers in BLUE.}
\end{center}

\vspace{0.75em}

\QBlock{Q1}{A research rover is located at the origin of a coordinate system and detects a beacon at position $\vec r = 8\hat i - 15\hat j$, where distances are measured in meters. Which of the following statements is correct?}{%
\Wrong{The distance between the beacon and the rover is $23\,\text{m}$.}{This adds the magnitudes of the components, but distance from the origin is the magnitude of the vector, not $|x|+|y|$. You must use the Pythagorean relation for perpendicular components.}
\Correct{The distance between the beacon and the rover is $17\,\text{m}$.}{The distance from the origin is the magnitude of the position vector: $|\vec r|=\sqrt{8^2+(-15)^2}$. This gives $|\vec r|=\sqrt{64+225}=\sqrt{289}=17\,\text{m}$, so this statement is correct.}
\Wrong{The distance between the beacon and the rover is $15\,\text{m}$.}{This confuses the magnitude of the $y$-component with the full distance from the origin. The total distance must include both perpendicular components.}
\Wrong{The distance between the beacon and the rover is $8\,\text{m}$.}{This uses only the $x$-component and ignores the vertical displacement. A single component is not the same as the vector's magnitude.}
}{B}

\QBlock{Q2}{If the displacement of an object is given by $\vec s(t)=\big((2.2\,\text{m/s}^3)t^3+(0.8\,\text{m/s}^2)t^2\big)\hat i$, which of the following statements is correct?}{%
\Wrong{This is a constant-acceleration motion.}{Because the displacement contains a cubic term in $t$, taking two derivatives produces an acceleration that depends on time. A time-dependent acceleration is not constant acceleration.}
\Wrong{The magnitude of the acceleration is $0.8\,\text{m/s}^2$.}{This mistakes a coefficient in the displacement function for the acceleration itself. Acceleration must be found from the second derivative of position with respect to time.}
\Correct{The magnitude of the acceleration at $t=0$ is $1.6\,\text{m/s}^2$.}{First differentiate displacement to get velocity: $\vec v(t)=\frac{d\vec s}{dt}=\big((6.6)t^2+(1.6)t\big)\hat i$. Differentiating again gives $\vec a(t)=\frac{d\vec v}{dt}=\big(13.2t+1.6\big)\hat i$, so at $t=0$, $|\vec a(0)|=1.6\,\text{m/s}^2$.}
\Wrong{The magnitude of the acceleration at $t=0$ is $0\,\text{m/s}^2$.}{This can result from confusing initial velocity with initial acceleration or from differentiating incorrectly. Evaluating the actual acceleration function at $t=0$ gives a nonzero value.}
}{C}

\QBlock{Q3}{A particle takes $5.0\,\text{s}$ to move from position $A$, where $\vec r_A = 6\hat i - 3\hat j$, to position $B$, where $\vec r_B = -2\hat i + 9\hat j$. Distances are measured in meters. Which of the following statements is correct?}{%
\Wrong{The distance traveled is $14.4\,\text{m}$.}{$\sqrt{(-2-6)^2+(9-(-3))^2}=\sqrt{208}\approx 14.4\,\text{m}$ is the magnitude of the displacement, not necessarily the distance traveled. The actual distance depends on the path taken between $A$ and $B$.}
\Wrong{The speed at $t=0$ must be $2.9\,\text{m/s}$.}{This treats displacement magnitude divided by elapsed time as the initial speed. That calculation gives only an average based on net displacement, not the instantaneous speed at the start.}
\Wrong{The acceleration at $t=0$ must be $0.58\,\text{m/s}^2$.}{There is not enough information to determine the instantaneous acceleration at $t=0$. Knowing only initial position, final position, and total time does not uniquely fix the motion.}
\Correct{None of the above is correct.}{The displacement is $\Delta \vec r = (-2-6)\hat i + (9-(-3))\hat j = -8\hat i + 12\hat j$, whose magnitude is $\sqrt{208}\approx 14.4\,\text{m}$. However, that is not necessarily the distance traveled, and neither the initial speed nor the initial acceleration is uniquely determined from only the endpoints and elapsed time, so none of the listed statements is correct.}
}{D}

\end{document}
